\section{Channel Reliability}

In the context of cooperative driving, the central problem is to guarantee the controller an information flow that is regular enough to support safe decisions. Network metrics such as Packet Delivery Ratio (PDR), latency, and jitter therefore acquire a direct operational meaning: their degradation does not remain confined to the communication layer, but is immediately reflected in the platoon's ability to maintain short gaps and absorb disturbances without amplifying them \cite{segata2015communication,giordano2019joint}.

The literature shows that no single radio technology can provide, in every scenario, the reliability required by safety-critical applications such as cooperative platooning \cite{segata2023multi-technology}. IEEE~802.11p is effective for short-range direct communications, but when many vehicles transmit on the same radio channel the nodes must contend for medium access, and message delivery can become less regular, with higher delays and packet losses; Visible Light Communication (VLC) can offer very efficient intra-vehicular links, but it depends on optical visibility and alignment; cellular technologies extend coverage and service continuity, but they introduce infrastructure dependence, link variability, and transients due to handovers \cite{ishihara2015improving,hardes2019communication,bazzi2020wireless}. This is where the interest in multi-technology architectures arises, combining interfaces with different operational weaknesses in order to improve the overall robustness of the system.

The work of Segata \emph{et al.} \cite{segata2023multi-technology} fits exactly into this line of research, proposing a system that monitors the quality of different interfaces, estimates degradation through PDR, and activates safe transitions between controllers and gap configurations through the \emph{SafeSwitch} logic. This contribution is the starting point of the present thesis for two reasons. The first is scientific: the paper clearly highlights the importance of jointly studying networking, control, and vehicle dynamics. The second is methodological: the proposed architecture is modular enough to be reproduced and updated on a more recent toolchain.

Reproducing that work with updated versions of OMNeT++, Veins, and PLEXE, and with Simu5G instead of SimuLTE, is not equivalent to a simple software migration. The system's behavior can in fact be influenced by the different cellular network model, by the evolution of the frameworks, and by the integration choices made between components. The thesis therefore also serves as an exercise in experimental reproducibility and portability assessment: the goal is to understand whether the logic proposed in the paper remains valid in an updated simulation environment, with channel reliability as the focus of the analysis \cite{nardini2020simu5g,segata2023multi-technology}.
